\section{Round Trip II}
\label{sec:cses-round-trip-ii}

\problemstatement{Given a \emph{directed} graph of $n$ nodes and $m$ edges,
find a directed cycle and print its nodes in order, or \texttt{IMPOSSIBLE}.}

\begin{patternbox}
A directed cycle is found by \textbf{DFS with three colours}: white
(unvisited), grey (on the current recursion stack), black (fully
processed). An edge to a \emph{grey} node is a back edge -- it closes a
directed cycle.
\end{patternbox}

\subsection*{Grey means ``on the current path''}

Colour a node grey when DFS enters it and black when it finishes. During
DFS at $u$, if a neighbour $v$ is grey, then $v$ is an ancestor still on the
recursion stack, so the path $v\to\cdots\to u$ plus the edge $u\to v$ is a
directed cycle. Reconstruct it via parent pointers from $u$ back to $v$. A
neighbour that is black is already finished and off the path, so it does not
indicate a cycle (unlike the undirected case, direction matters).

\begin{ideabox}[Back Edge to a Grey Ancestor]
The colour of the target distinguishes a cycle-closing back edge (grey,
still on the stack) from a harmless cross/forward edge (black, done).
Tracking parents lets the grey-to-current path be printed as the cycle.
\end{ideabox}

\subsection*{Algorithm}

\begin{enumerate}
  \item DFS from each white node; mark grey on entry, black on exit,
        storing parents.
  \item On finding a grey neighbour $v$ from $u$: record $v$ (cycle start)
        and $u$ (cycle end), stop.
  \item Reconstruct $v\to\cdots\to u\to v$ and print; else
        \texttt{IMPOSSIBLE}.
\end{enumerate}

\subsection*{Dry run}

\dryrun{$1\!\to\!2$, $2\!\to\!3$, $3\!\to\!1$.}

\begin{itemize}
  \item DFS greys $1,2,3$; from $3$ the edge to grey $1$ is a back edge.
  \item Cycle $1\to2\to3\to1$.
\end{itemize}

\subsection*{C++ implementation}

\begin{lstlisting}
#include <bits/stdc++.h>
using namespace std;

int n, m;
vector<vector<int>> adj;
vector<int> color, parent;                       // 0 white, 1 grey, 2 black
int cs = -1, ce = -1;

bool dfs(int u) {
    color[u] = 1;
    for (int v : adj[u]) {
        if (color[v] == 0) {
            parent[v] = u;
            if (dfs(v)) return true;
        } else if (color[v] == 1) {              // back edge to grey ancestor
            cs = v; ce = u; return true;
        }
    }
    color[u] = 2;
    return false;
}

int main() {
    cin >> n >> m;
    adj.assign(n + 1, {});
    for (int i = 0; i < m; i++) {
        int a, b; cin >> a >> b;
        adj[a].push_back(b);
    }
    color.assign(n + 1, 0);
    parent.assign(n + 1, 0);

    for (int s = 1; s <= n; s++)
        if (color[s] == 0 && dfs(s)) break;

    if (cs == -1) { cout << "IMPOSSIBLE\n"; return 0; }
    vector<int> cyc;
    cyc.push_back(cs);
    for (int v = ce; v != cs; v = parent[v]) cyc.push_back(v);
    cyc.push_back(cs);
    reverse(cyc.begin(), cyc.end());
    cout << cyc.size() << "\n";
    for (int v : cyc) cout << v << " ";
    cout << "\n";
    return 0;
}
\end{lstlisting}

\begin{complexitybox}
\textbf{Time Complexity.} $O(n+m)$ -- one DFS. \textbf{Space Complexity.}
$O(n+m)$, recursion depth $O(n)$ (use an explicit stack for very large $n$).
\end{complexitybox}

\begin{takeawaybox}
Directed cycle detection is DFS with grey/black colours: a back edge points
to a grey (on-stack) ancestor. The colour distinction -- absent in the
undirected version -- is exactly what respects edge direction.
\end{takeawaybox}
